% Nouveau EGA V (cote 25): archive cote 25, site page 77 (PDF page 78).
% typescript. Draft compared: none.
% AI-transcribed by Claude Opus 5.5 (claude-opus-5-5, adaptive thinking, effort high), Message Batches API, batch ega25b (msgbatch_01A7n5YuXubUPNK1zsepi8RG), 2026-11-13; checked against both supplied scan parts; see CHECK_DECISIONS.tsv.
\scanpage{25/77}
\mspage{\hand{72}}

une sous-variété linéaire donnée C  de codimension n.  A cet effet,
on aura besoin du

\uline{Lemme} 13.3.  Avec les notations précédentes pour \var{\xout{xx}}{}{typewriter strike-out}{ega25_25_077-v0001} X,Y,Z  ($X\hand{\supset} Y\hand{\supset} Z$
sous schémas fermés de $P^r_k$ ,\var{\inserted[typed]{de dimensions n,n-1 et \hand{$\leqslant$} n-2}}{de dimensions n,n-1 et \hand{$\leqslant$} n-2}{insertion}{ega25_25_077-v0002} X-Z\hand{,} \var{\deleted{lisse de dim n en tous ses points}}{}{deletion}{ega25_25_077-v0003},
Y-Z lisse\hand{s} \var{\deleted{de dim n-1 en tous ses points;}}{}{deletion}{ega25_25_077-v0004}\hand{,} Z de dimension \hand{$\leqslant$}n-2), et quitte
au besoin (lorsque k est de car p$\hand{>}$0)  de remplacer l'immersion projecti-%
ve  i:X$\hand{\rightarrow}$P$^r$  par  un ``multiple'' quelconque $\hand{\phi}_n i$ ($n \hand{\geqslant} 2$)  comme dans le
n$^o$ 9,  il existe une sous-variété linéaire  $L_o$ de P$^r$ , de codimension
n-1, et ayant les vertus suivantes:

a)  $L_o \hand{\cap} Z = \hand{\emptyset}$ .

b)  $L_o \hand{\cap} X$  est lisse de dimension 1.

c)  $L_o \hand{\cap} Y$   est lisse de dimension  0.

(NB  k désigne un corps infini, pas la peine ici qu'il soit parfait).

Admettons  ce lemme  et montrons comment on en déduit  l'existence
d'un voisinage ouvert U de x, contenu dans X-Y, et admettant un morphis-%
me élémentaire U$\hand{\longrightarrow}$V.  \var{\deleted{En vertu du n$^o$ 7 appliqué à l'espace projectif}}{}{deletion}{ega25_25_077-v0005}
\var{\deleted{$L_o$ et au sous-schéma lisse $L_o\hand{\cap} Y$ de celui-ci,}}{}{deletion}{ega25_25_077-v0006} \var{\over{i}{I}}{I}{writing over}{ega25_25_077-v0007}l existe une sous-%
variété linéaire C de $L_o$ , de codimension n dans P$^r$ i.e. de codimension
1 dans $L_o$ , \var{\inserted{ne rencontrant pas l'ensemble fini \var{\deleted{\illegible{1}}}{}{deletion}{ega25_25_077-v0009} $(L_o\cap Y)\cup\{x\}$.}}{ne rencontrant pas l'ensemble fini \var{\deleted{\illegible{1}}}{}{deletion}{ega25_25_077-v0009} $(L_o\cap Y)\cup\{x\}$.}{insertion}{ega25_25_077-v0008} \var{\deleted{telle que C $\cap$ X ( = C $\cap$ ($L_o\cap$ X) ) soit lisse sur k et de}}{}{deletion}{ega25_25_077-v0010}
\var{\deleted{dimension 0; de plus, on peut exiger que}}{}{deletion}{ega25_25_077-v0011} \var{\deleted{l'on ait  $C \hand{\cap} Y = \hand{\emptyset}$}}{}{deletion}{ega25_25_077-v0012} \var{\inserted{et $x\notin C$,}}{et $x\notin C$,}{insertion}{ega25_25_077-v0013} \var{\deleted{(même}}{}{deletion}{ega25_25_077-v0014}
\var{\deleted{argument).}}{}{deletion}{ega25_25_077-v0015} Soit  T = X$\hand{\cap}$C, de sorte que T est un sous-schéma de X,
étale sur k, non vide,\var{\inserted{ne contenant pas x,}}{ne contenant pas x,}{insertion}{ega25_25_077-v0016} et \uline{disjoint} de Y.  Considérons d'autre part le
sous-schéma Q de Grass$^{n-1}$(P) correspondant aux sous-variétés linéaires
de P$^r$ contenant \unsure{C|E}, de sorte que  Q est un \var{\xout{espace schéma}}{}{typewriter strike-out}{ega25_25_077-v0017} espace projectif
de dimension n-1,   en particulier il est lisse sur k de dim. n-1.
Alors $L_o$ correspond à un point \hand{$\xi$}$_o$ de Q(k).  Considérons d'autre part
(avec les notations générales introduites par ailleurs)  l'image inverse
$\underline{X}^{n-1}_Q$ de $\underline{X}^{n-1}$  par  l'immersion Q$\hand{\longrightarrow}$ Grass$^{n-1}$ , et de même les images

\var{\margin{77}}{\placedmargin{77}}{marginal note placed}{ega25_25_077-v0018}
